%%%%%%%%%%%%%%%%%%%%%%% file template.tex %%%%%%%%%%%%%%%%%%%%%%%%%%%%%%%%%%%%%%
%
% This is a general template file for the LaTeX package SVJour3
% for Springer journals.          Springer Heidelberg 2006/03/15
%
% Copy it to a new file with a new name and use it as the basis
% for your article. Delete % signs as needed.
%
% This template includes a few options for different layouts and
% content for various journals. Please consult a previous issue of
% your journal as needed.
%
%%%%%%%%%%%%%%%%%%%%%%%%%%%%%%%%%%%%%%%%%%%%%%%%%%%%%%%%%%%%%%%%%%%%%%%%%%%%%%%%
%
\documentclass[smallcondensed]{svjour3}
%
\smartqed  % flush right qed marks, e.g. at end of proof
%
\usepackage{graphicx}
\usepackage{textcomp}
\usepackage{amsmath, amsfonts}
\usepackage{mathabx}
\DeclareMathOperator*{\argmin}{arg\,min}
\DeclareMathOperator*{\argmax}{arg\,max}
\DeclareMathOperator*{\conv}{conv}
% etc.
%
% please place your own definitions here and don't use \def but
% \newcommand{}{}
%
% Insert the name of "your journal" with
\journalname{Journal of Global Optimization}
%


\title{Convex approximations of the feasible region of the pooling problem}

\author{Dag Haugland}

%\authorrunning{Short form of author list} % if too long for running head

\institute{Dag Haugland \at
Department of Informatics, University of Bergen,\\
P.O.Box 7803, N-5020 Bergen, Norway.\\
\email{dag@ii.uib.no}}
\date{\today}
%\date{Received: date / Accepted: date}
% The correct dates will be entered by the editor

\begin{document}

\maketitle
\begin{abstract}
The pooling problem is an extension of the minimum cost flow problem defined on a directed graph with three layers of nodes,
where quality constraints are introduced at each terminal node.
Flow entering the network at the source nodes has a given quality, at the internal nodes (pools) the entering flow is blended,
and then sent to the terminal nodes where all entering flow streams are blended again.
The resulting flow quality at the terminals has to satisfy given bounds.
The objective is to find a cost-minimizing flow assignment that satisfies network capacities and the terminals' quality specifications.

Recently, it was proved that the pooling problem is NP-hard, and that the hardness persists when the network has a unique pool.
In contrast, instances with only one source or only one terminal can be formulated as compact linear programs, and thereby solved in polynomial time.
In this work, it is proved that the pooling problem remains NP-hard even if there is only one quality constraint at each terminal.
Further, it is proved that the NP-hardness also persists if the number of sources and the number of terminals are no more than two,
and it is proved that the problem remains hard if all in-degrees or all out-degrees are at most two.
Examples of special cases in which the problem is solvable by linear programming are also given.
Finally, some open problems, which need to be addressed in order to identify more closely
the borderlines between polynomially solvable and NP-hard variants of the pooling problem, are pointed out.
\keywords{Pooling Problem \and Network flow \and Computational complexity \and Polynomial reduction}
\end{abstract}


\section{Introduction}
In many industrial applications,
conventional network flow models, such as the minimum cost flow and the maximum flow problems,
need to be extended in order to have practical relevance.
An extension often seen in the process industry, 
is that the network sources supply flow of unequal qualities,
while flow satisfying quality requirements is demanded at the terminals (sinks).
Flow streams of unequal qualities are blended when they enter a node jointly,
such that the quality of the flow leaving the node is averaged over the entering qualities.

Qualities are typically defined as the proportions of various components in the flow,
for instance an undesired contaminant or a desired nutrition.
With qualities being such proportions, the formulas for computing quality updates become simple:
The quality of the flow leaving a node is the convex combination of the entering qualities,
where the coefficients are proportional to the flow values.
Despite this simplicity, early industrial applications of optimization revealed that,
compared with the pure network flow models, the extension
in question represents a significant increase in computational challenge.

Referring to the above extension as the \emph{pooling problem},
Haverly \cite{haverly78} recognized the difficulty of solving it in terms of linear programming (LP).
Haverly's work was motivated by problems related to optimal mixing of crude oils.
Later, the industrial relevance of the problem, notably in petroleum refining \cite{baker85,dewitt89},
the food industry \cite{kallrath05}, and in waste water processing \cite{galan98,misener10b},
has been acknowledged by many authors.
The pooling problem survey by Misener and Floudas \cite{misener09b} has a comprehensive list of references to work in this area.

Adding quality constraints to the traditional network flow models leads very naturally to a \emph{bilinear} model.
Early works on inexact solution methods \cite{baker85,dewitt89,haverly78} take advantage of bilinearity,
utilizing the advantage of facing an LP when the values of a subset of the variables are fixed.
The idea underlying \emph{successive linear programming} (SLP) is to solve this LP,
use the solution as the base point of another linearization,
and to repeat this procedure until a fixpoint is reached.
By means of small numerical examples,
Haverly \cite{haverly78,haverly79} demonstrates that the solutions produced by different SLP-algorithms depend heavily on their initial linearization point.

Floudas and Visweswaran \cite{floudas90b} gave the first exact solution algorithm for the pooling problem.
Lower bounds (in a minimization problem) are computed by solving a Lagrangian relaxation, and upper bounds are computed by solving a projection of the original problem.
By proving that the two bounds converge, convergence to the optimal solution is also proved.
Lagrangian relaxation is also used in various global optimization algorithms \cite{adhya99,almutairi09,bental94} subsequent to \cite{floudas90b},
whereas the branch-and-bound algorithm by Foulds et al.\ \cite{foulds92} is based on convex relaxations of each bilinear term \cite{alkhayyal83,mccormick76}. 
A running time analysis of all these exact methods will reveal exponential worst-case behavior.
The above exact algorithms have later been improved \cite{audet04,dey15,misener11,sahinidis05,visweswaran96b},
and new and innovative methods, such as the integer programming technique by Dey and Gupte \cite{dey15}, are emerging.
Although these are capable of computing the global optimum of larger instances than the earliest algorithms could solve,
their exponential running times still confine their applicability to instances of moderate size.

Based on practical experience, a general conception of the pooling problem as inapproachable by LP was early established.
This has recently been confirmed in terms of a proof of NP-hardness \cite{alfaki13a}, and that the hardness persists even in network instances with only one internal node (pool).
Despite such a proof, many questions concerning the complexity of the problem are left open.
The goal of the current work is to enhance the understanding of the intractability of the pooling problem further.
To this end, the complexity result in \cite{alfaki13a} is complemented by several new theorems.
As the main contributions, it is proved in the current text that the problem remains NP-hard even if there is only one quality parameter (contaminant),
and even if there are only two sources and two sinks.
Further, NP-hardness of the problem in the case of networks where all nodes have in-degree at most two is proved.
By a similar proof, it is also demonstrated that the problem remains NP-hard in the case where all nodes have at most two out-neighbors.
Thus, several of the questions left open in \cite{dey15} are answered in the current work.

The remainder of the text is organized as follows:
In Section \ref{sec:pre}, the pooling problem is defined in rigorous terms, necessary mathematical notation is introduced,
and some preliminary results are given.


\section{Preliminaries} \label{sec:pre}

\subsection{Definitions} \label{sec:def}
Let $D=(S,P,T,A)$ be a directed graph,
where the disjoint node sets $S$, $P$, and $T$ consist of \emph{sources}, \emph{pools} and \emph{terminals}, respectively,
and the arc set $A\subseteq (S\times P)\cup(P\times T)$ links sources with pools and pools with terminals.
Let $K$ be a finite set, and refer to its elements as \emph{quality parameters}.
In the sequel, components of vectors and matrices defined over the sets $S$, $P$, $T$, and $K$,
are identified by a subscript if they correspond to $S$, $P$, or $T$,
and by a superscript if they correspond to $K$.
Define the vectors of unit cost $c\in\mathbb{R}^A$,
node capacity $b\in\mathbb{R}_+^{S\cup P\cup T}$, quality $q\in\mathbb{R}^{S\times K}$, and quality bounds $u\in\mathbb{R}^{T\times K}$.
For all sources $s\in S$ and quality parameters $k\in K$, $q_s^k$ denotes the quality, as defined by $k$, of the flow entering $D$ at $s$,
and for all terminals $t\in T$, $u_t^k$ denotes the upper bound on the quality (smaller values indicate better quality) of the flow leaving $D$ at $t$.

For each pool $p\in P$, the sets of neighbor sources and terminals are denoted by
$S_p=\left\{s\in S: (s,p)\in A\right\}$ and
$T_p=\left\{t\in T: (p,t)\in A\right\}$, respectively.
For each source $s\in S$ and each terminal $t\in T$, respectively,
the sets of neighbor pools are denoted by
$P_s=\left\{p\in P: (s,p)\in A\right\}$, and
$P_t=\left\{p\in P: (p,t)\in A\right\}$.
The arc set $A$ has a partition $\left(A_S,A_T\right)$, where $A_S=A\cap\left(S\times P\right)$ and $A_T=A\cap\left(P\times T\right)$
denote the sets of arcs incident to some source and to some terminal, respectively.

An $x\in\mathbb{R}_+^A$ is said to be a \emph{flow} in $D$ if it satisfies the capacity and flow conservation constraints:
\begin{eqnarray}
           & \sum_{p\in P_s}x_{sp}\leq b_s & s\in S, \label{eq:caps} \\
           & \sum_{s\in S_p}x_{sp}\leq b_p & p\in P, \label{eq:capi} \\
           & \sum_{p\in P_t}x_{pt}\leq b_t & t\in T, \label{eq:capt} \\
           & \sum_{s\in S_p}x_{sp}-\sum_{t\in T_p}x_{pt} = 0~~~~ & p\in P. \label{eq:flowcons}
\end{eqnarray}

\noindent
Let a flow polytope be defined as $F(D,b) = \left\{x\in\mathbb{R}_+^A: \mbox{(\ref{eq:caps})--(\ref{eq:flowcons}) are satisfied}\right\}$.
For each $x\in F(D,b)$, define $w\in\mathbb{R}^{P\times K}$ to be a \emph{quality matrix} of $x$ if
$w_p^k=\sum_{s\in S_p}q_s^kx_{sp}/\sum_{t\in T_p}x_{pt}$ for all $k\in K$ and for all $p\in P$ for which
$\sum_{t\in T_p}x_{pt}>0$.
If $\sum_{t\in T_p}x_{pt}=0$, $w_p^k$ can take an arbitrary value.
That is, the quality vector $\left(w_p^k\right)_{k\in K}$ of the flow blended at pool $p$, is a convex combination of the source qualities $\left(q_s^k\right)_{k\in K}$,
with coefficients $x_{sp}/\sum_{t\in T_p}x_{pt}$ ($s\in S_p$).
This definition of $w$ is referred to as the \emph{linear blending} assumption.
At a terminal $t\in T$ where $\sum_{p\in P_t}x_{pt}>0$, the upper quality bounds $u_t^k$ apply,
and linear blending at $t$ yields $\sum_{p\in P_t}w_p^kx_{pt}/\sum_{p\in P_t}x_{pt}\leq u_t^k$.
Therefore, $x$ is said to be a \emph{feasible} flow in $D$ if $\sum_{p\in P_t}w_p^kx_{pt}\leq u_t^k\sum_{p\in P_t}x_{pt}$ ($t\in T$, $k\in K$)
for any quality matrix $w$ of $x$.
This leads to our problem definition:

\begin{problem} \label{prob:p} The \textsc{Pooling Problem}:
Find a feasible flow $x$ in $D$ minimizing $\sum_{(i,j)\in A}c_{ij}x_{ij}$.
\end{problem}

The assumption that there are no direct arcs from sources to terminals can be made without loss of generality,
because any such arc can be replaced by a new pool along with two arcs connecting the pool to the source and the terminal, respectively.

Lower quality bounds, that is, constraints on the form $\sum_{p\in P_t}w_p^kx_{pt} \geq \ell_t^k\sum_{p\in P_t}x_{pt}$, where $\ell_t^k$ are given constants for all $t\in T$ and $k\in K$,
do not represent an extension of (\ref{eq:obj})--(\ref{eq:posflow}).
For all quality parameters $k\in K$ subject to some lower bound, it is possible to extend $K$ by a new quality parameter $k'$,
to let $q_s^{k'}=-q_s^k$ for all $s\in S$, and to define the upper bounds $u_t^{k'}=-\ell_t^k$ for all $t\in T$.
Whenever relevant, instances of the problem where (\ref{eq:qualbnd}) is replaced by
$\ell_t^k\sum_{p\in P_t}x_{pt}\leq\sum_{p\in P_t}w_p^kx_{pt}\leq u_t^k\sum_{p\in P_t}x_{pt}$ are therefore referred to as instances of the \textsc{Pooling Problem}.
When the number of quality parameters is an issue, however, all quality parameters subject also to lower bounds have to be counted twice.


\subsection{Preliminary results} \label{sec:res}

The definition of Problem \ref{prob:p} gives directly a bilinear formulation in flow and quality variables \cite{haverly78},
often referred to as the P-formulation:

\begin{eqnarray}
\min_{x,w} & \sum_{(i,j)\in A}c_{ij}x_{ij}, & \label{eq:obj} \\
           & \sum_{s\in S_p}q_s^kx_{sp}-w_p^k\sum_{t\in T_p}x_{pt} = 0~~~~ & p\in P, k\in K, \label{eq:qualcons} \\
           & \sum_{p\in P_t}w_p^kx_{pt}-u_t^k\sum_{p\in P_t}x_{pt} \leq 0~~~~ & t\in T, k\in K, \label{eq:qualbnd} \\
           & x\in F(D,b). \label{eq:posflow}
\end{eqnarray}

\noindent
It is readily observed that for a pool $p\in P$ where $\sum_{t\in T_p}x_{pt}=0$ (a terminal $t\in T$ where $\sum_{p\in P_t}x_{pt}=0$),
constraints (\ref{eq:qualcons}) (constraints (\ref{eq:qualbnd})) are satisfied regardless of the value of $w$.

For any $x,y\in\mathbb{R}^n$, let $[x,y]$ denote the line segment with end points $x$ and $y$.

\begin{definition} \label{def:sets}
Let $\Omega\subseteq\mathbb{R}^n$ be a compact set.
We say that $x\in\Omega$ \emph{sees} $y\in\Omega$ via $\Omega$ if $[x,y]\subseteq\Omega$.
The \emph{convex kernel} of $\Omega$ is the set $\ker\Omega$ consisting of the points $x\in\Omega$ that sees all $y\in\Omega$.
\end{definition}

A set $C\subseteq\Omega$ is called a \emph{maximal convex subset} of $\Omega$ if $C$ is convex and
$\conv\left(C\cup\{y\}\right)\setminus\Omega\neq\emptyset$ for all $y\in\Omega\setminus C$.

\begin{lemma} \label{lem:all}
For all $x\in\Omega$, there exists some maximal convex subset $C$ of $\Omega$ such that $x\in\Omega$.
\end{lemma}
\begin{proof}
Because $\{x\}$ is a convex subset of $\Omega$, there exists a sequence $\left\{C_i\right\}_{i=1}^{\infty}$ of convex sets such that
$x\in C_i\subseteq C_{i+1}\subseteq\Omega$ for all $i=1,2,\ldots$.
Since also $\bigcup_{i=1}^{\infty}C_i$ is a convex subset of $\Omega$ containing $x$, the result follows from Zorn's lemma.
\qed
\end{proof}

The following result is stated without proof in \cite{kenelly69}:
\begin{proposition} \label{prop:equiv}
$\ker\Omega$ is the intersection of all maximal convex subsets of $\Omega$.
\end{proposition}
\begin{proof}
Assume $x\in\ker\Omega$, and that there exists a maximal convex subset $C$ of $\Omega$ such that $x\not\in C$.
Because $x$ sees all $y\in C$, we have $[x,y]\subseteq\Omega$.
It follows that $\conv\left(C\cup\{x\}\right)\subseteq\Omega$, contradicting maximality of $C$.

Assume next that $x\in C$ for all maximal convex subsets $C$ of $\Omega$, and consider any $y\in\Omega$.
By Lemma \ref{lem:all}, $y\in C'$ for some maximal convex subset $C'$ of $\Omega$.
Hence, $[x,y]\subseteq C'\subseteq\Omega$, and $x$ sees $y$ via $\Omega$.
\qed
\end{proof}

Observe that Proposition \ref{prop:equiv} includes the cases where $\Omega$ has an empty convex kernel.
Sets with non-empty kernels are usually referred to as \emph{star-shaped} sets.

\begin{theorem} \label{thm:kerpp}
Let $\Omega$ be the set of feasible flows in $D$.
Then $\ker\Omega$ is the set of flows satisfying the following homogeneous linear inequalities:
\begin{eqnarray}
           & \sum_{s\in S_p}\left(q_s^k-u_t^k\right)x_{sp}\leq 0~~~~ & p\in P, t\in T[p], k\in K, \label{eq:kerqual} \\
           & x\in F(D,b), \label{eq:kerflow}
\end{eqnarray}
where, for all $p\in P$, $T[p]\subseteq T_p$ is the set of terminals $t$
for which there exists a $w_p\in\conv\left\{q_s:s\in S_p\right\}$ such that $w_p\leq u_t$.
\end{theorem}
\begin{proof}
Let $x,y\in\Omega$, and assume $x$ satisfies (\ref{eq:kerqual})--(\ref{eq:kerflow}).
By proving that $[x,y]\subseteq\Omega$, we prove that $x\in\ker\Omega$.
\qed
\end{proof}

\begin{thebibliography}{99}
\bibitem{adhya99}
Adhya, N., Tawarmalani, M., and Sahinidis, N.V.
\newblock A Lagrangian approach to the pooling problem.
\newblock \emph{Industrial \& Engineering Chemical Research},
\newblock \textbf{38}(5), pp.\ 1965--1972, 1999.

\bibitem{alfaki13a}
Alfaki, M. and Haugland, D.
\newblock Strong formulations for the pooling problem.
\newblock \emph{Journal of Global Optimization},
\newblock \textbf{56}(3), pp.\ 897--916, 2013.

\bibitem{alkhayyal83}
Al-Khayyal, F.A. and Falk, J.E.
\newblock Jointly constrained biconvex programming. 
\newblock \emph{Mathematics of Operations Research},
\newblock \textbf{8}(2), pp.\ 273--286, 1983.

\bibitem{almutairi09}
Almutairi, H. and Elhedhli, S.
\newblock A new Lagrangian approach to the pooling problem.
\newblock \emph{Journal of Global Optimization},
\newblock \textbf{45}(2), pp.\ 237--257, 2009.

\bibitem{audet04}
Audet, C., Brimberg, J., Hansen, P., Le Digabel, S., and Mladenovi\'{c}, N.
\newblock Pooling problem: alternate formulations and solution methods. 
\newblock \emph{Management Science},
\newblock \textbf{50}(6), pp.\ 761--776, 2004.

\bibitem{baker85}
Baker, T.E. and Lasdon, L.S.
\newblock Successive linear programming at Exxon.
\newblock \emph{Management Science},
\newblock \textbf{31}(3), pp.\ 264--274, 1985.

\bibitem{bental94}
Ben-Tal, A., Eiger, G., and Gershovitz, V.
\newblock Global minimization by reducing the duality gap.
\newblock \emph{Mathematical Programming},
\newblock \textbf{63}(1--3), pp.\ 193--212, 1994.

\bibitem{dey15}
Dey, S. and Gupte, A.
\newblock Analysis of MILP techniques for the pooling problem.
\newblock \emph{Operations Research}, accepted up to minor revision.

\bibitem{dewitt89}
DeWitt, C.W., Lasdon, L.S., Waren, A.D., Brenner, D.A., and Melham, S.
\newblock OMEGA: An improved gasoline blending system for Texaco. 
\newblock \emph{Interfaces},
\newblock \textbf{19}(1), pp.\ 85--101, 1989.

\bibitem{floudas90b}
Floudas, C.A. and Visweswaran, V.
\newblock A global optimization algorithm (GOP) for certain classes of nonconvex NLPs: I. Theory. 
\newblock \emph{Computers \& Chemical Engineering}, 
\newblock \textbf{14}(12), pp.\ 1397--1417, 1990.

\bibitem{foulds92}
Foulds, L.R., Haugland, D., and J\"{o}rnsten, K.
\newblock A bilinear approach to the pooling problem.
\newblock \emph{Optimization},
\newblock \textbf{24}, pp.\ 165--180, 1992.

\bibitem{galan98}
Galan, B. and Grossmann, I.E.
\newblock Optimal design of distributed wastewater treatment networks.
\newblock \emph{Industrial \& Engineering Chemistry Research}, 
\newblock \textbf{37}(10), pp.\ 4036--4048, 1998.

\bibitem{garey79}
Garey, M.R. and Johnson, D.S.
\newblock \emph{Computers and Intractability: A Guide to the Theory of NP-completeness}.
\newblock W.H.Freeman, New York, New York, 1979.

\bibitem{garey76}
Garey, M.R., Johnson, D.S., and Stockmeyer L.
\newblock Some simplified NP-complete graph problems.
\newblock \emph{Theoretical Computer Science},
\newblock \textbf{1}, pp.\ 237--267, 1976.

\bibitem{haverly78}
Haverly, C.A.
\newblock Studies of the behavior of recursion for the pooling problem.
\newblock \emph{ACM SIGMAP Bulletin},
\newblock \textbf{25}, pp.\ 19--28, 1978.

\bibitem{haverly79}
Haverly, C.A.
\newblock Behavior of recursion models - more studies.
\newblock \emph{ACM SIGMAP Bulletin},
\newblock \textbf{26}, pp.\ 22--28, 1979.

\bibitem{kallrath05}
Kallrath, J.
\newblock Solving planning and design problems in the process industry using mixed integer and global optimization.
\newblock \emph{Annals of Operations Research},
\newblock \textbf{140}(1), pp.\ 339--373, 2005.

\bibitem{kenelly69}
Kenelly, J.W., Hare Jr., W.R., Evans, B.D., and Ludesche, W.H.
\newblock Convex components, extreme points, and the convex kernel
\newblock \emph{Proceedings of the American Mathematical Society}
\newblock \textbf{21}(1), pp.\ 83--87, 1969.

\bibitem{kohli94}
Kohli, R., Krishnamurti, R., and Mirchandani, P.
\newblock The minimum satisfiability problem.
\newblock \emph{Discrete Mathematics},
\newblock \textbf{7}, pp.\ 275--283, 1994.

\bibitem{misener09b}
Misener, R. and Floudas, C.A.
\newblock Advances for the pooling problem: modeling, global optimization, and computational studies.
\newblock \emph{Applied Computational Mathematics},
\newblock \textbf{8}, pp.\ 3--22, 2009.

\bibitem{misener10b}
Misener, R. and Floudas, C.A.
\newblock Global optimization of large-scale generalized pooling problems: Quadratically constrained MINLP models.
\newblock \emph{Industrial \& Engineering Chemistry Research}, 
\newblock \textbf{49}, pp.\ 5424--5438, 2010.

\bibitem{misener11}
Misener, R., Thompson, J.P, and Floudas, C.A.
\newblock APOGEE: Global optimization of standard, generalized, and extended pooling
problems via linear and logarithmic partitioning schemes.
\newblock \emph{Computers and Chemical Engineering},
\newblock \textbf{35}, pp.\ 876--892, 2011.

\bibitem{mccormick76}
McCormick, G.P.
\newblock Computability of global solutions to factorable nonconvex programs: Part 1-convex underestimating problems.
\newblock \emph{Mathematical Programming},
\newblock \textbf{10}(1), pp.\ 147--175, 1976.

\bibitem{sahinidis05}
Sahinidis, N.V. and Tawarmalani, M.
\newblock Accelerating branch-and-bound through a modeling language construct for relaxation-specific constraints.
\newblock \emph{Journal of Global Optimization},
\newblock \textbf{32}(2), pp.\ 259--280, 2005.

\bibitem{visweswaran96b}
Visweswaran, V. and Floudas, C.A.
\newblock Computational results for an efficient implementation of the GOP algorithm and its variants.
\newblock In Grossmann, I.E.(ed.): \emph{Global Optimization in Chemical Engineering},
\newblock pp.\ 111–153. Kluwer Academic Publishers, 1996.

\end{thebibliography}


\end{document}
